\begin{table}[H]
\centering
\begin{tabular}{p{3cm} p{7cm} p{6cm}}
\toprule
\textbf{\small Estimand} & \textbf{What It Measures} & \textbf{Example Question} \\
\midrule
\textbf{\footnotesize Marginal Means (MM)} 
& The average probability that a profile with a given attribute level (e.g., ``Male'') is chosen across all possible matches, including ties. It can be interpreted as proportion of binary choices in favour of given attribute level
& ``How often does a `Male' profile get selected, on average, when compared to all other profiles?'' \\[1em]

\midrule

\textbf{\footnotesize Average Marginal Component Effects (AMCE)} 
& On average, how does shifting from one attribute level (e.g., ``Female'') to another (e.g., ``Male''), under the same distribution, effect the probability of the entire choice set being chosen.
& ``If we changed a profile’s gender from female to male (holding everything else constant), by how much would the likelihood of being chosen increase/decrease on average?'' \\[1em]

\midrule

\textbf{\footnotesize Average Component Preferences (ACP)} 
& It measure the direct within-pair preference for one level over another, ignoring pairs where the two profiles share that same attribute level. This allows to compare within-pair preference with other attributes. 
& ``When respondents see one candidate who is male and another who is female, which one do they pick more often?'' \\
\bottomrule
\end{tabular}
\caption{Measures and Research Questions for MM, AMCE, and ACP}
\label{tab:estimands}
\end{table}